% ===========================================================================
% Vue logique en composants.
%
% Relevé au 2026-09-22 : 18 modules d'endpoints exposant 173 routes sous
% /api/v1, 33 services applicatifs, 16 modules de modèles (25 tables),
% 3 clients d'intégration. Le découpage en couches ci-dessous est celui du
% dépôt (backend/app/), pas une reconstruction de convenance.
% ===========================================================================
\begin{tikzpicture}[node distance=5mm]

  % ─── Couche client ──────────────────────────────────────────────────────
  \node[pqbox, text width=31mm] (spa)   {Front-office\\{\scriptsize candidat, visiteur}};
  \node[pqbox, right=of spa, text width=31mm] (bo) {Back-office\\{\scriptsize recruteur, RH, admin}};
  \node[pqbox, right=of bo, text width=31mm] (pub) {Pages publiques\\{\scriptsize liens à jeton}};

  \node[pqbox, below=7mm of bo, text width=101mm]
    (coreng) {Noyau Angular — services HTTP, intercepteur de jeton,
              gardes de route, modèles typés};

  \begin{scope}[on background layer]
    \node[pqgroup, fit=(spa)(pub)(coreng)] (clientbox) {};
  \end{scope}
  \node[pqcaption, anchor=south west, xshift=2mm] at (clientbox.north west)
    {Client — Angular 19.2, composants autonomes, chargement différé};

  % ─── Frontière HTTP ─────────────────────────────────────────────────────
  \node[pqbox, below=9mm of coreng, text width=101mm, fill=white]
    (nginx) {nginx 1.27 — service des fichiers statiques, relais \texttt{/api/},
             relais WebSocket, en-têtes de sécurité};

  % ─── Couche API ─────────────────────────────────────────────────────────
  \node[pqcore, below=9mm of nginx, text width=101mm]
    (api) {Couche API — FastAPI, 18 modules d'endpoints, 173 routes,
           2 points WebSocket};

  \node[pqbox, below=6mm of api.south west, anchor=north west, text width=32mm]
    (auth) {Authentification\\{\scriptsize JWT, 2FA, OAuth}};
  \node[pqbox, right=of auth, text width=32mm]
    (rbac) {Autorisation\\{\scriptsize Casbin, refus par défaut}};
  \node[pqbox, right=of rbac, text width=32mm]
    (valid) {Validation\\{\scriptsize schémas Pydantic}};

  \begin{scope}[on background layer]
    \node[pqboundary, fit=(api)(auth)(valid)] (apibox) {};
  \end{scope}

  % ─── Couche services ────────────────────────────────────────────────────
  % 15 mm sous la couche API : il faut la place du titre de groupe, qui se
  % pose au-dessus du cadre en pointillés.
  \node[pqbox, below=15mm of auth, text width=32mm]
    (srec) {Recrutement\\{\scriptsize offres, candidatures, CV}};
  \node[pqbox, right=of srec, text width=32mm]
    (sia) {Traitement du langage\\{\scriptsize analyse, notation, sémantique}};
  \node[pqbox, right=of sia, text width=32mm]
    (sproc) {Processus\\{\scriptsize entretiens, tests, contrats}};

  \node[pqbox, below=of srec, text width=32mm]
    (sml) {Apprentissage\\{\scriptsize prédiction salariale}};
  \node[pqbox, below=of sia, text width=32mm]
    (srh) {RH interne\\{\scriptsize employés, budget, carrière}};
  \node[pqbox, below=of sproc, text width=32mm]
    (srep) {Décision et diffusion\\{\scriptsize rapports, notifications}};

  \begin{scope}[on background layer]
    \node[pqgroup, fit=(srec)(srep)] (svcbox) {};
  \end{scope}
  \node[pqcaption, anchor=south west, xshift=2mm] at (svcbox.north west)
    {Services applicatifs — 33 modules, règles métier};

  % ─── Couche persistance ─────────────────────────────────────────────────
  \node[pqbox, below=8mm of srh, text width=101mm]
    (orm) {Accès aux données — SQLAlchemy 2.0 asynchrone, 25 tables,
           schéma piloté par Alembic (18 révisions)};

  \node[pqstore, below=7mm of orm, xshift=-32mm] (pg)
    {PostgreSQL 16\\+ pgvector};
  \node[pqstore, right=14mm of pg] (rd) {Redis 7\\{\scriptsize cache, session}};
  \node[pqstore, right=14mm of rd] (fs) {Volume CV\\{\scriptsize fichiers déposés}};

  % ─── Intégrations ───────────────────────────────────────────────────────
  % Hors du cadre de l'API : ce sont des sorties vers l'extérieur, pas une
  % sous-couche. Décalé vers le bas pour ne pas heurter le trait du cadre.
  \node[pqext, right=13mm of api.east, yshift=-6mm, text width=26mm, anchor=west]
    (integ) {Intégrations\\{\scriptsize 3 clients,\\ toutes facultatives}};

  % ─── Liens verticaux ────────────────────────────────────────────────────
  \draw[pqflow] (coreng.south) -- (nginx.north);
  \draw[pqflow] (nginx.south)  -- (api.north);
  \draw[pqflow] (auth.south)  -- (srec.north);
  \draw[pqflow] (rbac.south)  -- (sia.north);
  \draw[pqflow] (valid.south) -- (sproc.north);
  \draw[pqflow] (sml.south)   -- (sml.south |- orm.north);
  \draw[pqflow] (srh.south)   -- (orm.north);
  \draw[pqflow] (srep.south)  -- (srep.south |- orm.north);
  \draw[pqflow] (orm.south -| pg.north) -- (pg.north);
  \draw[pqflow] (orm.south -| rd.north) -- (rd.north);
  \draw[pqflow] (orm.south -| fs.north) -- (fs.north);
  \draw[pqopt]  (api.east) -- (integ.west);

\end{tikzpicture}
